% =============================================================================
\section{Discretização por invariância ao impulso}\label{sec:invimp}
% =============================================================================

\subsection*{Teoria mínima}

\paragraph{Definição.} Dado um protótipo analógico $H_a(s)$ com resposta ao
impulso $h_a(t)$, o filtro digital é definido por
\[
  \boxed{\ h[n]=T_s\,h_a(nT_s)\ }
\]
(a multiplicação por $T_s$ é o que faz $|H(e^{j\omega})|\approx|H_a(j\Omega)|$
na banda-base). No plano $z$, cada polo $p_k$ do protótipo vira
\[
  z_k=e^{p_kT_s}\qquad\text{e}\qquad
  H(z)=\sum_k\frac{T_sA_k}{1-e^{p_kT_s}z^{-1}},
  \qquad A_k=\text{resíduo de }H_a \text{ em }p_k .
\]

\paragraph{Caso de polos complexos (o do Desafio 08).} Para
\[
  H_a(s)=\frac{s}{s^2+2\zeta\omega_ns+\omega_n^2},\qquad
  \sigma=\zeta\omega_n,\ \omega_d=\omega_n\sqrt{1-\zeta^2},
\]
\[
  h_a(t)=e^{-\sigma t}\Bigl[\cos\omega_dt-\frac{\sigma}{\omega_d}
  \operatorname{sen}\omega_dt\Bigr]u(t),
\]
e, escrevendo $r=e^{-\sigma T_s}$ e $\theta=\omega_dT_s$,
\[
  \boxed{\ H(z)=\frac{T_s\left[1-r\left(\cos\theta+
  \frac{\sigma}{\omega_d}\operatorname{sen}\theta\right)z^{-1}\right]}
  {1-2r\cos\theta\,z^{-1}+r^{2}z^{-2}}\ }
\]
de modo que
\[
  b_0=T_s,\quad
  b_1=-T_sr\Bigl(\cos\theta+\frac{\sigma}{\omega_d}\operatorname{sen}\theta\Bigr),
  \quad a_1=-2r\cos\theta,\quad a_2=r^{2}.
\]

\paragraph{O que o método preserva e o que ele destrói.}
\begin{itemize}
  \item \textbf{Preserva:} a resposta ao impulso amostrada (por construção) e
        as posições dos \textbf{polos} ($z=e^{pT_s}$), logo
        estabilidade e frequência de ressonância.
  \item \textbf{Destrói/Inventa:} os \textbf{zeros}. Zeros \emph{não} são
        mapeados por $z=e^{sT_s}$: o numerador de $H(z)$ é apenas combinado
        pelos resíduos. Um zero em $s=0$ (passa-altas) simplesmente
        desaparece.
  \item \textbf{Aliasing:} a resposta em frequência é a soma de réplicas
        (Capítulo~\ref{sec:amostragem}):
        \[
          H(e^{j\omega})=\sum_kH_a\!\left(j\frac{\omega-2\pi k}{T_s}\right)
          +\frac{T_sh_a(0^+)}{2}.
        \]
  \item \textbf{Conclusão prática:} o método é adequado para filtros
        \emph{limitados em banda} (passa-baixas e passa-faixa com
        $|H_a(j\Omega)|\to0$ em alta frequência). Ele \textbf{falha} para
        passa-altas e rejeita-faixa, onde $|H_a(j\Omega)|\to\text{const}\neq0$.
\end{itemize}

\ideia{Teste rápido de validade: se $|H_a(j\Omega)|$ não tende a zero quando
$\Omega\to\infty$, o somatório de réplicas diverge e o filtro digital obtido
não terá nada a ver com o protótipo. Nesse caso usa-se a transformação
bilinear (Capítulo~\ref{sec:bilinear}).}

\subsection*{Exercícios resolvidos}

% ---------------------------------------------------------------- 15.1
\exer{Discretize $h_a(t)=e^{-1000t}u(t)$ com $T_s=0{,}1$ ms. Dê $H(z)$, o
ganho em DC e compare com o analógico $H_a(0)=1/1000$.}

\textbf{Solução.} Polos: $p=-1000$ (real) $\Rightarrow z_0=e^{-1000\cdot10^{-4}}
=e^{-0{,}1}=0{,}9048374180$, e o resíduo é 1, logo
\[
  H(z)=\frac{T_s}{1-e^{-1000T_s}z^{-1}}
      =\frac{10^{-4}}{1-0{,}9048374180\,z^{-1}} .
\]
Ganho em DC:
\[
  H(1)=\frac{10^{-4}}{1-0{,}9048374180}=1{,}0508332\times10^{-3}
  \quad\text{contra}\quad H_a(0)=\frac{1}{1000}=1{,}000\times10^{-3}.
\]
Erro de $+5{,}08\%$ — resíduo da aproximação ``soma de Riemann'' da
integração usada no mapeamento (o erro cresce com $aT_s$; aqui
$aT_s=0{,}1$).

\resp{$H(z)=\dfrac{10^{-4}}{1-0{,}9048374180z^{-1}}$;
$H(1)=1{,}0508332\times10^{-3}$ ($+5{,}08\%$ em relação a $1/a$).}

% ---------------------------------------------------------------- 15.2
\exer{\textbf{(Item 2 do Desafio 08)} Aplique o método ao protótipo
$H_a(s)=\dfrac{s}{s^2+200\pi s+(12000\pi)^2}$ com $T_s=25\ \mu$s e escreva os
coeficientes de $H(z)$.}

\textbf{Solução.} Do Capítulo~\ref{sec:laplace}: $\sigma=100\pi$,
$\omega_d=100\pi\sqrt{14399}$. Então
\[
  r=e^{-\sigma T_s}=e^{-100\pi\cdot25\times10^{-6}}
   =0{,}992176780293,
\]
\[
  \theta=\omega_dT_s=100\pi\sqrt{14399}\cdot25\times10^{-6}
   =0{,}942445070585\ \text{rad}=0{,}2999895832\pi\ \text{rad/amostra}.
\]
Com $T_s=2{,}5\times10^{-5}$:
\[
  \begin{array}{rcl}
    b_0&=&T_s=2{,}500000000000\times10^{-5},\\[1mm]
    b_1&=&-T_sr\left(\cos\theta+\frac{\sigma}{\omega_d}\operatorname{sen}\theta\right)
        =-1{,}474755715063\times10^{-5},\\[1mm]
    a_1&=&-2r\cos\theta=-1{,}166426294332,\\[1mm]
    a_2&=&r^{2}=0{,}984414763352 .
  \end{array}
\]
Polos: $z=0{,}992176780e^{\pm j0{,}29999\pi}$ ($|z|<1$: estável). Note que
$\theta\approx0{,}3\pi$, que é justamente o $\omega_0$ do Item 3.

\resp{$b_0=T_s$, $b_1=-1{,}4747557\times10^{-5}$, $a_1=-1{,}1664263$,
$a_2=0{,}9844148$; polos em $0{,}992177e^{\pm j0{,}3\pi}$.}

% ---------------------------------------------------------------- 15.3
\exer{\textbf{(desafio de reflexão)} Mostre que o método degenera para o
passa-altas $H_a(s)=\dfrac{s}{s+2000\pi}$ com $T_s=25\ \mu$s, calculando
$|H(1)|$ e $|H(-1)|$ e explicando a causa.}

\textbf{Solução.} Anti-transformando:
\[
  h_a(t)=\delta(t)-2000\pi e^{-2000\pi t}u(t),
\]
que tem o termo impulsivo $\delta(t)$ — a marca registrada dos passa-altas.
Amostrando,
\[
  h[n]=T_s\delta[n]-T_s\,2000\pi\,\alpha^{n}u[n],\qquad
  \alpha=e^{-2000\pi T_s}=0{,}8546359992,
\]
\[
  H(z)=T_s-\frac{T_s\cdot2000\pi}{1-\alpha z^{-1}}
      =\frac{(T_s-T_s\,2000\pi)-T_s\alpha z^{-1}}{1-\alpha z^{-1}} .
\]
Ganhos nos extremos da banda digital:
\[
  |H(1)|=\left|\frac{T_s-T_s\cdot2000\pi-T_s\alpha}{1-\alpha}\right|
  =1{,}080570 \quad(+0{,}6731\ \text{dB}),
\]
\[
  |H(-1)|=\left|\frac{(T_s-T_s\cdot2000\pi)+T_s\alpha}{1+\alpha}\right|
  =0{,}084671 \quad(-21{,}4453\ \text{dB}) .
\]
O ganho \emph{cai} de $1{,}08$ em DC para $0{,}085$ em Nyquist: o
``passa-altas'' virou \textbf{passa-baixas}. Causas (três argumentos
independentes):
\begin{enumerate}
  \item $|H_a(j\Omega)|=\frac{\Omega}{\sqrt{\Omega^2+(2000\pi)^2}}\to1$
        quando $\Omega\to\infty$: cada réplica da identidade de Poisson tem
        ganho $\approx1$ e o somatório \textbf{não converge} — nenhum $F_s$
        conserta isso.
  \item O zero em $s=0$, que definia o passa-altas, \textbf{não é mapeado}
        para $z=1$: o zero do filtro digital ficou em
        $z=-\hat{b}_1/\hat{b}_0\approx-1{,}36\times10^{-4}$ (praticamente na
        origem).
  \item Algebricamente, $H(z)$ é a soma de um termo constante
        ($T_s\delta[n]$) com um \textbf{passa-baixas de 1ª ordem} (o polo em
        $\alpha$); a diferença de uma constante com um passa-baixas é uma
        função \emph{decrescente} da frequência.
\end{enumerate}
Solução correta para passa-altas: \textbf{transformação bilinear com
pré-deformação} (Capítulo~\ref{sec:bilinear}).

\resp{$|H(1)|=1{,}080570$ e $|H(-1)|=0{,}084694$: filtro passa-baixas;
falha porque $|H_a|\not\to0$ em HF (réplicas não decaem) e porque zeros não
mapeiam ($s=0\nrightarrow z=1$).}
